\documentclass[11pt]{article}
\usepackage{amsmath,amsthm,amssymb}
\usepackage[margin=1in]{geometry}
\usepackage[colorlinks=true,linkcolor=blue,urlcolor=blue,citecolor=blue]{hyperref}
\usepackage{mathtools}
\usepackage{enumitem}
\usepackage{booktabs}
\usepackage[round,authoryear]{natbib}

\theoremstyle{plain}
\newtheorem{theorem}{Theorem}
\newtheorem{lemma}{Lemma}
\newtheorem{corollary}{Corollary}
\newtheorem{proposition}{Proposition}
\theoremstyle{definition}
\newtheorem{assumption}{Assumption}
\newtheorem{definition}{Definition}
\theoremstyle{remark}
\newtheorem{remark}{Remark}

\newcommand{\E}{\mathbb{E}}
\newcommand{\PP}{\mathbb{P}}
\newcommand{\Pn}{\mathbb{P}_n}
\newcommand{\R}{\mathbb{R}}
\newcommand{\cX}{\mathcal{X}}
\newcommand{\cT}{\mathcal{T}}
\newcommand{\cS}{\mathcal{S}}
\newcommand{\cP}{\mathcal{P}}
\newcommand{\ez}{e_0}
\newcommand{\mz}{\mu_0}
\newcommand{\wz}{w_0}
\newcommand{\ehat}{\widehat{e}}
\newcommand{\mhat}{\widehat{\mu}}
\newcommand{\what}{\widehat{w}}
\newcommand{\that}{\widehat{\tau}}
\newcommand{\thetahat}{\widehat{\theta}}
\newcommand{\pihat}{\widehat{\pi}}
\newcommand{\Vhat}{\widehat{V}}
\newcommand{\Dw}{D_w}
\newcommand{\Dm}{D_\mu}
\newcommand{\Ltwo}[1]{\lVert #1 \rVert_{2}}
\newcommand{\Lw}[1]{\lVert #1 \rVert_{2,w}}
\DeclareMathOperator*{\argmin}{arg\,min}
\DeclareMathOperator{\Var}{Var}
\DeclarePairedDelimiter{\abs}{\lvert}{\rvert}
\DeclarePairedDelimiter{\card}{\lvert}{\rvert}

\title{Interpretable Causal Inference with Optimal Decision Trees\\
\large Notation, assumptions, and main results}
\author{}
\date{2026-08-06 --- revision 2, after adversarial check}

\begin{document}
\maketitle

\begin{abstract}
Notation, assumptions and main-result statements for the paper, without proofs, so
the symbol system and assumption set can be checked before drafting.
\textbf{Revision~2} repairs four blocking defects found in revision~1; these and
all lesser corrections are listed in Section~\ref{sec:changes}. The largest
structural change: Part~III is now stated on the \emph{realised} errors of the
fitted nuisances rather than on approximation errors, which makes the bias
constant exactly $1/\pi$, removes a hidden per-nuisance rate requirement, and
makes the allowance estimable rather than assumed.
\end{abstract}

\tableofcontents

\section{Notation}\label{sec:notation}

\subsection{Data, target, score}

\begin{center}\begin{tabular}{@{}ll@{\hspace{1.5em}}l@{}}
\toprule
Symbol & Meaning & Notes \\ \midrule
$O=(X,A,Y)$ & one observation & $O_1,\dots,O_n$ i.i.d.\ $\sim P$ \\
$X\in\cX$, $d$ & covariates, dimension & $d$, \textbf{not} $p$ \\
$A\in\{0,1\}$, $Y\in\R$ & treatment, outcome & $Y(a)$ potential outcome \\
$\Pn$ & empirical measure & $\Pn f=n^{-1}\sum_if(O_i)$ \\
$\pi=\PP(A=1)$ & treated share & $\pihat=\Pn A$ \\
$m_1(x)$ & $\E[Y\mid A=1,X=x]$ & appears only in $V$ \\
$\sigma^2_a(x)$ & $\Var(Y\mid A=a,X=x)$ & \\
$\theta_0$ & ATT $=\E[Y(1)-Y(0)\mid A=1]$ & target \\
$\psi$, $V$, $\Vhat$ & EIF, its variance, plug-in & \eqref{eq:score},
   \eqref{eq:vdecomp}, \eqref{eq:vhat} \\
$z_{1-\alpha/2}$ & normal quantile & \\
\bottomrule
\end{tabular}\end{center}

\subsection{Covariate space}

\begin{center}\begin{tabular}{@{}ll@{\hspace{1.5em}}l@{}}
\toprule
Symbol & Meaning & Notes \\ \midrule
$\cX=\prod_{j\le d}\cX_j$ & finite product space & \\
$J$ & levels per coordinate & \textbf{not} $K$ \\
$M=\card\cX\le J^d$ & number of cells & \emph{fixed} \\
$p_x=\PP(X=x)$ & cell probability & $p_{\min}=\min\{p_x>0\}$ \\
$s$ & \emph{local} sparsity & active coordinates per region \\
$\bar\alpha$ & harmonic-mean smoothness & continuous tier only \\
$K$, $n_k$ & folds, fold size & Part~I only \\
\bottomrule
\end{tabular}\end{center}

Since $s$ is local, the active set may differ across regions and $\gamma_0$ may
depend on all $d$ coordinates globally while $s$ stays small
(Remark~\ref{rem:additive}).

\subsection{Nuisances and norms}

$\ez(x)=\PP(A=1\mid X=x)$; $\mz(x)=\E[Y\mid A=0,X=x]$ (controls only);
$\gamma_{0,j}$ for $j\in\{e,\mu\}$; $\eta=(\ehat,\mhat)$;
$\wz=\ez/(1-\ez)$ the ATT weight; $c$ the positivity constant
(A\ref{ass:causal}); $\rho(\cdot,\cdot)$ the loss (\textbf{not} $L$);
$\Ltwo{f}^2=\E f(X)^2$ and $\Lw{f}^2=\E[\{1-\ez(X)\}f(X)^2]$.

\begin{remark}[Norm equivalence: the constant, stated correctly]\label{rem:normconst}
Since $c\le1-\ez\le1$,
\[
  c\,\Ltwo f^2\le\Lw f^2\le\Ltwo f^2,\qquad
  \sqrt c\,\Ltwo f\le\Lw f\le\Ltwo f .
\]
The constant is $1/c$ for \emph{squared} norms and $c^{-1/2}$ for \emph{norms}.
Revision~1 wrote $1/c$ for norms, and that slip was the proximate cause of a wrong
constant in the honest interval; both are fixed here.
\end{remark}

\subsection{Trees, partitions, selection}

$\tau$ a recursive binary axis-aligned tree partition, $\ell$ a leaf,
$\card\tau$ its leaf count; $\bar L$ the (fixed) leaf budget;
$\cT_{\bar L}=\{\tau:\card\tau\le\bar L\}$, a fixed finite set;
$\cS_j$ the sufficient class (Def.~\ref{def:sufficient}); $\that_j$ the selected
partition; $m_n$ the minimum leaf count (\textbf{not} fold size); $\lambda_n$ the
leaf penalty; $R^{(j)},R^{(j)}_n$ population and empirical risk; $\Delta_j$ the
margin (Lemma~\ref{lem:margin}); $\Delta_{\min}$ its uniformity floor
(A\ref{ass:margin-floor}). Selection is
\begin{equation}\label{eq:select}
  \that_j\in\argmin_{\tau\in\cT_{\bar L}}
  \bigl\{R^{(j)}_n(\tau)+\lambda_n\card\tau\bigr\},
\end{equation}
a \emph{global} minimum; leaf values are within-leaf means (treated fraction for
$\ehat$, control mean for $\mhat$).

\subsection{Approximation versus realised error}

$\Dw=\Lw{\what-\wz}$ and $\Dm=\Lw{\mhat-\mz}$ are the \emph{realised} errors
(Def.~\ref{def:realised}), with debiased estimators $\widehat D_j$ and one-sided
upper bounds $U_w,U_\mu$ (Prop.~\ref{prop:estimable}). $\delta_j$ is the
\emph{approximation} error (Def.~\ref{def:delta}), with decay exponent $q_j$.
$\delta_{\mathrm{fid}}=\thetahat-\thetahat_{\mathrm{cf}}$ is the fidelity
diagnostic (Lemma~\ref{lem:spectest}).

\begin{remark}[Three different objects]\label{rem:D-vs-delta}
$D_j$ is a property of the \emph{fitted} nuisance on the sample in hand;
$\delta_j$ an approximation-theoretic minimum over $\cT_{\bar L}$; they satisfy
$D_j=\delta_j+O_p(\sqrt{\bar L/n})$. \textbf{Revision~1 conflated them}, which
made the honest interval false (Remark~\ref{rem:whyD}). $\delta_{\mathrm{fid}}$ is
a difference of \emph{estimators} and bounds neither.
\end{remark}

\subsection{Symbol collisions resolved}\label{sec:collisions}

\begin{center}\begin{tabular}{@{}lp{5.8cm}l@{}}
\toprule
Symbol & Competing meanings & Resolution \\ \midrule
$s$ & sparsity; smoothness; leaf count &
  $s$; $\bar\alpha$; $L$ \\
$K$ & folds; levels per coordinate & $K$; $J$ \\
$p$ & dimension; cell probability; weak-$\ell_p$ & $d$; $p_x$; weak-$\ell_r$ \\
$m$ & minimum leaf count; fold size & $m_n$; $n_k$; penalty $\lambda_n$ \\
$L$ & leaf count; loss function & $L$; $\rho$ \\
$\delta$ & approximation error; realised error; fidelity diagnostic &
  $\delta_j$; $D_j$; $\delta_{\mathrm{fid}}$ \\
\bottomrule
\end{tabular}\end{center}

\section{The estimator}

The efficient influence function for the ATT \citep{hahn1998} is
\begin{equation}\label{eq:score}
  \psi(O;\theta,\eta)=\frac1\pi\Bigl[A\{Y-\mu(X)-\theta\}
   -(1-A)\,w(X)\{Y-\mu(X)\}\Bigr],\qquad w=\frac e{1-e}.
\end{equation}
Solving $\Pn\psi(\cdot;\theta,\widehat\eta)=0$,
\begin{equation}\label{eq:att-estimator}
  \thetahat=\frac1{n\pihat}\sum_{i=1}^n\Bigl[A_i\{Y_i-\mhat(X_i)\}
   -(1-A_i)\,\what(X_i)\{Y_i-\mhat(X_i)\}\Bigr].
\end{equation}

\begin{remark}[Sign, and why $\pihat$ costs nothing]\label{rem:sign}
The control-arm term carries a \textbf{minus}; an earlier draft displayed a plus
while defining the score as \eqref{eq:score}, so the displayed estimator was not
the one the score defines. Also, $\pihat$ is not substituted --- it \emph{emerges}
as $n^{-1}\sum A_i$ from the coefficient of $\theta$. With $g:=\pi\psi$,
$\thetahat-\theta_0=\Pn g(\cdot;\theta_0,\widehat\eta)/\pihat$ exactly, so
\[
 \thetahat-\theta_0=\Pn\psi(\cdot;\theta_0,\widehat\eta)
  -\frac{\pihat-\pi}{\pi}\Pn\psi(\cdot;\theta_0,\widehat\eta)+O_p(n^{-3/2}),
\]
whose second term is $O_p(n^{-1})$. No \emph{further} correction for $\pihat$ is
needed: the $-\theta_0A/\pi$ term inside \eqref{eq:score} already is the
normalisation adjustment. (Revision~1's display omitted a factor $1/\pi$ and a
sign; the order was right.)
\end{remark}

The second-order remainder is exactly bilinear:
\begin{equation}\label{eq:bilinear}
  \E[\psi(O;\theta_0,\eta)]
  =\frac1\pi\E\Bigl[\{1-\ez(X)\}\{\wz(X)-w(X)\}\{\mz(X)-\mu(X)\}\Bigr].
\end{equation}
This is the textbook doubly-robust computation
\citep{robins1995,vanderlaan2003,kennedy2024} and should be cited, not claimed.

\begin{lemma}[Exact bias bound, constant $1/\pi$]\label{lem:biasbound}
$\bigl\lvert\E[\psi(O;\theta_0,\widehat\eta)]\bigr\rvert\le\pi^{-1}\Dw\Dm$.
\end{lemma}

\noindent Cauchy--Schwarz with respect to $(1-\ez)\,dP$ splits the weight as
$\sqrt{1-\ez}$ onto \emph{both} factors; no power of $c$ appears. Converting to
the conventional unweighted propensity error costs $c^{-3/2}$:
\begin{equation}\label{eq:convert}
  \bigl\lvert\E[\psi(O;\theta_0,\widehat\eta)]\bigr\rvert
  \le\frac1{\pi c^{3/2}}\Ltwo{\ehat-\ez}\,\Lw{\mhat-\mz},
\end{equation}
using $(1-\ez)\abs{\wz-w}=\abs{\ez-e}/(1-e)\le c^{-1}\abs{\ez-e}$ and then
Remark~\ref{rem:normconst}. Form \eqref{eq:convert} is what Part~I's rate
condition uses; Lemma~\ref{lem:biasbound} is what Part~III uses.

\section{Assumptions}

\begin{assumption}[Identification and one-sided positivity]\label{ass:causal}
Consistency; $\{Y(0),Y(1)\}\perp A\mid X$; $\pi>0$; and there is $c\in(0,1)$ with
$1-\ez(x)\ge c$ for every $x\in\cX$.
\end{assumption}

\begin{remark}[Positivity is one-sided]\label{rem:positivity}
\textbf{Weakened from earlier drafts}, which assumed $\ez\in[c,1-c]$. Keep two
statements apart. \emph{Identification} needs only $\ez(x)<1$ pointwise on
$\{\ez>0\}$, with no uniform constant: $\E[Y(1)\mid A=1]$ comes from treated units
with no positivity, and cells with $\ez(x)=0$ carry weight $0$ in
$\theta_0=\pi^{-1}\E[\ez(m_1-\mz)]$. \emph{Inference} needs the uniform
$1-\ez\ge c$ and nothing more: $\wz$ is bounded iff $1-\ez\ge c$;
\eqref{eq:bilinear} involves only $1-\ez,1-\ehat$; $V$ is finite whenever
$\E[\ez^2\sigma_0^2/(1-\ez)]<\infty$, for which $1-\ez\ge c$ plus
A\ref{ass:moments} suffices; and the margin needs control mass
$p_x\{1-\ez(x)\}\ge p_xc$. No lower bound on $\ez$ is used anywhere, and
$\ez(x)=0$ exactly is harmless ($\wz=0$, $\ehat_\ell=0$ with zero variance, $\mz$
well defined).

Two refinements. The margin \emph{splits}: $\Delta_e$ needs \emph{no} positivity
(its excess risk is $\Ltwo{\Pi_\tau\ez-\ez}^2>0$ once a leaf mixes two positive-
probability cells with different $\ez$), while $\Delta_\mu$ needs $c$ because the
$\mu$-risk is control-weighted. And the cross-entropy loss--norm link, the one
place two-sidedness previously entered, has a constant diverging in the
\emph{fitted} values, which are under our control --- clipping candidates
two-sidedly removes any need to \emph{assume} $\ez\ge c$. Under squared-error
fitting the issue does not arise (A\ref{ass:lossnorm}).
\end{remark}

\begin{assumption}[Second moments and non-degeneracy]\label{ass:moments}
$\E[Y^2\mid A=a,X=x]\le C_Y<\infty$ for all $a,x$; and $V>0$.
\end{assumption}

\noindent Revision~1 conditioned on $X$ only, while the quantity used is
conditional on $(A,X)$ (the bridge $\E[Y^2\mid A=0,X]\le C_Y/c$ is available, but
stating it directly is cleaner); and $V>0$ was nowhere assumed although the Wald
interval divides by $\sqrt{\Vhat}$.

\begin{assumption}[Loss--norm link]\label{ass:lossnorm}
Nuisances are fitted by squared error, for which the excess risk \emph{equals}
$\Ltwo{f-\gamma_0}^2$, so $C_{\mathrm{ln}}=1$.
\end{assumption}

\noindent For binary $A$, $\E[(A-f)^2\mid X]=\ez(1-\ez)+(\ez-f)^2$, so the excess
risk is $\Ltwo{f-\ez}^2$ exactly, for any $\ez\in[0,1]$, with no positivity.
Earlier drafts said $C_{\mathrm{ln}}=\tfrac12$. Two scope notes: the within-leaf
treated fraction minimises \emph{both} squared error and log-loss over constants,
so the \textbf{leaf refit is loss-invariant} and the loss affects only the
structure search \eqref{eq:select}; and under A\ref{ass:sparsity} the best tree
has excess risk exactly zero, so \emph{no} norm-to-risk conversion is needed in
Part~II at all --- it is live only in Part~III and Part~I's oracle inequality.

\begin{assumption}[Global optimisation]\label{ass:global}
$\that_j$ solves \eqref{eq:select} globally over $\cT_{\bar L}$, $\lambda_n\to0$.
\end{assumption}

\begin{assumption}[Clipping at the positivity constant]\label{ass:clip}
Fitted propensities are clipped so $1-\ehat(x)\ge c$, with $c$ \textbf{the same
constant as A\ref{ass:causal}}.
\end{assumption}

\begin{remark}[Why clip, why only above, why the level matters]\label{rem:clip}
A leaf of only treated units gives $\ehat=1$ and a weight dividing by zero;
probability at most $(1-c)^m$ for a leaf of size $m$ --- small but not zero, and
it will occur across many replications. Clipping only from \emph{above} is
strictly better than two-sided: since A\ref{ass:causal} puts $\ez\in[0,1-c]$,
clipping is the Euclidean projection onto that interval, hence non-expansive,
$\abs{\mathrm{clip}(\ehat)-\ez}\le\abs{\ehat-\ez}$, so every $L_2$ rate can only
improve; a two-sided clip loses this wherever $\ez<c$.

\textbf{The clip level must equal $c$.} If the software clips at $c'>c$ then on
cells with $\ez>1-c'$ we get $\what\to w'\ne\wz$: $\thetahat$ stays consistent by
double robustness, but the influence function changes, so $V$,
Corollary~\ref{cor:variance} and every coverage claim break.

Clipping does \emph{not} disturb Lemma~\ref{lem:cross}: that lemma rests on
$\sum_{i\in\ell,A_i=0}(Y_i-\mhat_\ell)=0$, a property of the \emph{outcome} leaf
means, which clipping never touches, and clipping is a deterministic map of
$\ehat$ so the coefficients stay $\sigma(X_{1:n},A_{1:n})$-measurable. In
particular $m_n\to\infty$ is \textbf{not} required for this purpose --- a claim
revision~1's check initially made and which is withdrawn.
\end{remark}

\begin{assumption}[Minimum leaf mass, counted in the right arm]\label{ass:mass}
Every leaf of $\that_e$ holds $\ge m_n$ observations; every leaf of $\that_\mu$
holds $\ge m_n$ \textbf{control} observations; $\bar L\,m_n\lesssim n$. On the
null event that a leaf of $\that_\mu$ has no controls, $\mhat$ takes the parent
leaf's control mean.
\end{assumption}

\begin{remark}[Why controls, and why balance]\label{rem:mass}
\textbf{Corrected from revision~1}, which counted observations: $\mhat_\ell$ is
the \emph{control} mean, so an all-treated leaf makes it $0/0$, and clipping
$\ehat$ does nothing for it. For balance, a control leaf mean has variance
$\sigma^2_\ell/\{np_\ell(1-\bar e_\ell)\}$, so
\[
 \E\Lw{\mhat-\mz}^2\asymp\frac1n\sum_\ell
   \frac{p_\ell(1-\bar e_\ell)\sigma^2_\ell}{p_\ell(1-\bar e_\ell)}
  =\frac1n\sum_\ell\sigma^2_\ell\lesssim\frac{\bar L}n,
\]
the control factor \textbf{cancelling exactly in the weighted norm} --- an
independent reason Definition~\ref{def:realised} uses $\Lw\cdot$; the unweighted
version carries a $1/c$. Balanced leaves, $\min_\ell p_\ell\gtrsim1/\bar L$, are
what make this $\bar L/n$; on a finite space with $p_{\min}>0$ this follows from
A\ref{ass:causal} automatically.
\end{remark}

\begin{assumption}[Leaf budget: fixed]\label{ass:budget}
$\bar L$ is a fixed constant, $\card{\that_j}\le\bar L$, and $\cT_{\bar L}$ is a
fixed finite set.
\end{assumption}

\begin{remark}[Why fixed, and what it buys]\label{rem:fixed-tier}
\textbf{Changed from revision~1}, which permitted a growing $L_n=o(\sqrt n)$
alongside a fixed finite $\cX$ --- inconsistent, since $M$ fixed forces
$\bar L\le M=O(1)$. With $\bar L,M,d,J$ fixed:
$\log\card{\cT_{\bar L}}\asymp\bar L\log(dJ)=O(1)$, so the union bound costs a
\emph{constant}; no exponential tail bound is used anywhere; A\ref{ass:moments}
suffices with no sub-Gaussian condition; and $\bar L\log n=o(\sqrt n)$ is vacuous.
\textbf{Part~II must therefore not be motivated as escaping an
$\exp(C\bar L\log n)$ entropy cost} --- our own finite-space assumption forbids
that cost from existing.

The self-influence reading survives: one observation's influence on its own leaf
mean is $O(1/n_{\mathrm{leaf}})=O(\bar L/n)$, which is $o(n^{-1/2})$ whenever
$\bar L=o(\sqrt n)$ --- automatic here. This is \citet{chen2022}'s leave-one-out
stability condition specialised to leaf means; ours is the verification, not the
idea. The growing regime is scoped out (Section~\ref{sec:open}).
\end{remark}

\begin{definition}[Sufficient class, defined intrinsically]\label{def:sufficient}
$\cS_j=\{\tau\in\cT_{\bar L}:\Pi_\tau\gamma_{0,j}=\gamma_{0,j}\ P\text{-a.s.}\}$,
where $\Pi_\tau$ projects onto functions constant on $\tau$ (for $j=\mu$, in
$\Lw\cdot$).
\end{definition}

\begin{remark}[The counterexample that forces the intrinsic definition]
\label{rem:pinwheel}
\textbf{This replaces revision~1's ``$\tau$ refines $\tau^*_j$'', a blocking
error.} Revision~1 took $\tau^*_j$ minimal with $\gamma_{0,j}$ constant on it and
asserted $\bar L\ge\card{\tau^*_j}$ is ``equivalently $\delta_j=0$''. False: a
minimal \emph{box} partition need not be a \emph{tree} partition.

Let $d=2$, $\cX=\{0,1,2\}^2$, all nine cells of positive probability, and let
$\gamma_0$ be constant on and distinct across
\[
 A=\{x_1{=}0,x_2\in\{0,1\}\},\ B=\{x_1\in\{0,1\},x_2{=}2\},\
 C=\{x_1{=}2,x_2\in\{1,2\}\},
\]
\[
 D=\{x_1\in\{1,2\},x_2{=}0\},\ E=\{(1,1)\}.
\]
These five boxes partition $\cX$, so $\card{\tau^*}=5$; but no first split is
clean ($x_1$ cut $\{0\}\mid\{1,2\}$ splits $B$, $\{0,1\}\mid\{2\}$ splits $D$,
symmetrically on $x_2$), and the minimal \emph{tree} refinement needs six leaves.
So at $\bar L=5=\card{\tau^*}$: $\delta_j>0$ strictly, $\cS_j=\emptyset$,
Lemma~\ref{lem:selection}'s conclusion false (that probability is $0$), and
Theorem~\ref{thm:main} with no valid route. Definition~\ref{def:sufficient}
sidesteps box-versus-tree and the support qualifier at once.
\end{remark}

\begin{assumption}[Structural sparsity]\label{ass:sparsity}
$\cS_e\ne\emptyset$ and $\cS_\mu\ne\emptyset$; equivalently
$\delta_e=\delta_\mu=0$.
\end{assumption}

\begin{remark}[Sparsity, not correct specification]\label{rem:not-exact}
On a finite space \emph{every} function is piecewise constant --- the singleton
partition always works and \emph{is} tree-realisable, via $M-1$ splits. So
A\ref{ass:sparsity} asserts nothing about functional form; its entire content is
that a \emph{small} tree suffices, i.e.\ that the nuisances are structurally
sparse and hence displayable. It must not be called exactness or correct
specification, and it is not vulnerable to ``you assumed the truth is a step
function''.
\end{remark}

\begin{remark}[Why trees rather than a sparse additive model]\label{rem:additive}
The minimal tree size is at most $J^s$ for $s$ active coordinates (equal to $2^s$
only when $J=2$; revision~1 wrote $2^s$ unconditionally), often far less:
\emph{hierarchical} dependence, where the active set differs by region, costs as
few as $s+1$ leaves. A\ref{ass:sparsity} is expensive precisely for
\emph{additive} structure ($x_1+x_2+x_3$ needs a leaf per configuration, trees'
worst case) and cheap precisely for region-varying active sets. That asymmetry,
not flexibility, is the reason to use trees. A concrete example must accompany
this remark.
\end{remark}

\begin{assumption}[Margin floor, for uniform statements only]\label{ass:margin-floor}
$\min_j\Delta_j\ge\Delta_{\min}>0$.
\end{assumption}

\begin{remark}[Pointwise versus uniform coverage]\label{rem:uniformity}
Lemma~\ref{lem:selection} gives
$\PP(\that_j\notin\cS_j)\le\card{\cT_{\bar L}}e^{-c_1n\Delta_j^2}$, which is
$o(1)$ for fixed $P$ but \textbf{not uniformly} over
$\{P:\delta_e=\delta_\mu=0\}$, since $\Delta_j$ can be arbitrarily small there
(two cells with nearly equal nuisance values). This is the Leeb--P\"otscher
non-uniformity of post-selection intervals. Without A\ref{ass:margin-floor} all
coverage claims below are \emph{pointwise in $P$}; with it they hold uniformly
over $\cP(\Delta_{\min})=\{P:\delta_e=\delta_\mu=0,\min_j\Delta_j\ge\Delta_{\min}\}$,
at the cost of one assumption and no new proof. Part~III does not repair this: it
widens by a term that is exactly zero in this class.
\end{remark}

\section{Part I: cross-fitted inference with optimal-tree nuisances}

\begin{assumption}[Rate condition, product form]\label{ass:rate}
$\Ltwo{\ehat-\ez}\cdot\Lw{\mhat-\mz}=o_p(n^{-1/2})$; in the continuous tier,
$\bar\alpha_e/(2\bar\alpha_e+s_e)+\bar\alpha_\mu/(2\bar\alpha_\mu+s_\mu)>1/2$.
\end{assumption}

\begin{remark}[Product form, and why the mixed pairing is right]\label{rem:product}
The per-nuisance $o_p(n^{-1/4})$ requirement forces $\bar\alpha_j>s_j/2$; since
$s$ is an integer supremum of active-set cardinalities with no fractional
relaxation, and $\bar\alpha\le1$ for piecewise-constant trees, it forces
$\boldsymbol{s=1}$ --- locally univariate, excluding every smooth two-way
interaction. The product form needs only that the two ratios sum past $1/2$.

The pairing is \emph{unweighted} for $e$, \emph{weighted} for $\mu$, exactly as
\eqref{eq:convert} delivers. The reason is \emph{not} that the weight ``sits on
the $\mu$ factor'' --- symmetric Cauchy--Schwarz weights both
(Lemma~\ref{lem:biasbound}). It is that one full factor of $(1-\ez)$ is absorbed
by $(1-\ez)\abs{\wz-w}=\abs{\ez-e}/(1-e)$, leaving $\mu$ to be converted by
$c^{-1/2}$.
\end{remark}

\begin{theorem}[Cross-fitted ATT with optimal-tree nuisances]\label{thm:crossfit}
Under A\ref{ass:causal}--A\ref{ass:lossnorm}, A\ref{ass:rate}, and $K$-fold
cross-fitting with $K$ fixed,
$\sqrt n(\thetahat_{\mathrm{cf}}-\theta_0)\rightsquigarrow N(0,V)$.
\end{theorem}

\noindent \emph{Rung two of the ladder.} No leaf budget, no sparsity assumption,
and $K$ trees per nuisance rather than one --- individually readable and auditable
fold by fold, already a gain over a black-box nuisance, but not a single displayed
model. This is also the honest answer to ``how would you relax
A\ref{ass:sparsity}''.

\section{Part II: two trees, full sample, no splitting}

Structure \emph{and} leaf values from all $n$: no folds, no splitting, no
cross-fitting, no intersection, no averaging over fits.

\begin{lemma}[Sufficiency and invariance of the estimand]\label{lem:invariance}
Under A\ref{ass:sparsity}, for every $(\tau_e,\tau_\mu)\in\cS_e\times\cS_\mu$ the
population plug-in estimand equals $\theta_0$.
\end{lemma}

\noindent For $j=\mu$ the leaf value is the $p_x\{1-\ez(x)\}$-weighted average of
$\mz$ over the leaf, so this needs positive control mass per leaf --- supplied by
A\ref{ass:causal}. That is where one-sided positivity does work inside Part~II.

\begin{lemma}[The margin is free under global optimisation]\label{lem:margin}
Under A\ref{ass:causal}--A\ref{ass:sparsity},
$\Delta_j:=\min_{\tau\in\cT_{\bar L}\setminus\cS_j}
\{R^{(j)}(\tau)-\min_{\tau'}R^{(j)}(\tau')\}>0$.
\end{lemma}

\begin{remark}[Global optimisation supplies the margin; greedy does not]
\label{rem:greedy}
$\cT_{\bar L}$ is finite and any $\tau\notin\cS_j$ has a leaf containing two
positive-probability cells with distinct $\gamma_{0,j}$, so its excess risk is
strictly positive; a minimum of finitely many positives is positive. Under greedy
recursive splitting no margin exists: if $\ez$ depends on $X$ only through
$\mathbf1\{x_1=x_2\}$ the marginal risk reduction of \emph{either} correct first
split is exactly zero, so $\Delta=0$. The result therefore genuinely requires an
exact solver --- one place where optimal rather than CART-style trees do
mathematical rather than rhetorical work.
\end{remark}

\begin{lemma}[Selection lands in the sufficient class]\label{lem:selection}
Under A\ref{ass:causal}--A\ref{ass:global} and A\ref{ass:sparsity},
$\PP(\that_j\in\cS_j)\to1$; if $Y$ is bounded,
$\PP(\that_j\notin\cS_j)\le\card{\cT_{\bar L}}e^{-c_1n\Delta_j^2}$.
\end{lemma}

\noindent Only the \emph{exponential rate} needs a tail condition;
$\PP(\that_j\in\cS_j)\to1$ holds under A\ref{ass:moments} alone, since for fixed
$\tau$ the empirical risk converges by the law of large numbers and uniformity
over a fixed finite candidate set is free.

\subsection{The correct decomposition}

With $\varepsilon_i=Y_i-\mz(X_i)$, $\Delta\mu=\mhat-\mz$, $\Delta w=\what-\wz$,
and $\phi(O;\eta)=A\{Y-\mu(X)-\theta_0\}-(1-A)w(X)\{Y-\mu(X)\}$, we have
\emph{exactly}, with no expansion,
\begin{equation}\label{eq:skeleton}
  \pihat\,(\thetahat-\theta_0)=\Pn\phi(\eta_0)+U-T,
\end{equation}
\[
 U:=\Pn[h(O)\Delta\mu(X)],\quad h(O):=\wz(X)(1-A)-A,\qquad
 T:=\Pn[(1-A)\Delta w(X)\{Y-\mhat(X)\}],
\]
with $\E[h\mid X]=\wz(1-\ez)-\ez=0$ pointwise.

\begin{remark}[There is no separate remainder term that ``vanishes'']
\label{rem:no-vanish}
\textbf{Corrected from revision~1}, which said the bilinear remainder vanishes
under A\ref{ass:sparsity}. A\ref{ass:sparsity} sets the \emph{approximation} error
to zero; the second-order product is absorbed inside $T$, and what remains is
driven by \emph{estimation} error of the leaf values. Concretely
$\abs U=O_p(\bar L/n)$ --- negligible because A\ref{ass:budget} makes $\bar L$
fixed (and in a growing regime because $\bar L=o(\sqrt n)$) --- and
$\abs T=O_p(\sqrt{L_e}/n)$ by Lemma~\ref{lem:cross}. So A\ref{ass:sparsity} kills
the \emph{bias} component while A\ref{ass:mass}--A\ref{ass:budget} kill the
\emph{estimation} component. Note $U$, not $T$, is the binding constraint.
\end{remark}

\begin{lemma}[Cross term via leaf-wise conditional variance]\label{lem:cross}
Under A\ref{ass:causal}--A\ref{ass:sparsity}, with $\that_e\ne\that_\mu$ carried
separately, $\abs T=O_p(\sqrt{L_e}/n)=o_p(n^{-1/2})$.
\end{lemma}

\noindent \emph{Mechanism.} Since $\sum_{i\in\ell,A_i=0}(Y_i-\mhat_\ell)=0$,
subtracting any leaf-constant from the coefficient is free; take
$a_i:=\Delta w(X_i)-\overline{\Delta w}_{\ell(i)}$. On
$\{\that_\mu\in\cS_\mu\}$, $\mz$ is leaf-constant, so
$Y_i-\mhat_\ell=\varepsilon_i-(\mhat_\ell-\mu_{0,\ell})$ and the second piece is
annihilated, giving \emph{exactly} $T=n^{-1}\sum_{i:A_i=0}a_i\varepsilon_i$. For
fixed $\tau_\mu$ the $a_i$ are $\sigma(X_{1:n},A_{1:n})$-measurable, so
$\E[T\mid X,A]=0$ and Chebyshev applies; centring is a projection, so
$\sum a_i^2\le c^{-4}\sum_{\ell'}n_{0,\ell'}(\ehat_{\ell'}-\ez)^2=O_p(L_e)$ with
\emph{no} leaf-mass assumption.

\begin{remark}[No $\log n$; and Cauchy--Schwarz is fatal, not lossy]\label{rem:cs}
\textbf{Corrected from revision~1}, whose $O_p(\sqrt{\bar L}\log n/n)$ carried a
$\log n$ contradicting Remark~\ref{rem:fixed-tier}'s own $O(1)$ entropy. The
$\log n$ arises only from a Bernstein-plus-union argument needed when
$\cT_{\bar L}$ grows; with $M,d,J,\bar L$ fixed the maximum of finitely many
$O_p(\sqrt{L_e}/n)$ variables is $O_p(\sqrt{L_e}/n)$.

Cauchy--Schwarz instead gives a second factor converging to the control-arm
residual standard deviation $\E[(1-\ez)\Var(Y\mid A=0,X)]^{1/2}=O_p(1)$, bounded
\emph{away from zero} --- irreducible outcome noise, not $\delta_\mu$. The bound
is then $O_p(\sqrt{L_e/n})\ge O_p(n^{-1/2})$ for every $L_e\ge1$, so no leaf
budget rescues it: the failure is at the CLT scale. The structural reason is that
$T$ is \emph{not} of product form --- it pairs the propensity error with outcome
\emph{noise}, so the argument must exploit conditional mean-zero-ness and never
smallness. Cross terms therefore do not vanish under A\ref{ass:sparsity} either;
exactness only makes the $(Y-\mhat)\to\varepsilon$ substitution exact, which is
A\ref{ass:sparsity}'s only role in this lemma.
\end{remark}

\begin{lemma}[Uniform asymptotic linearity on the sufficient class]
\label{lem:uniform}
Under A\ref{ass:causal}--A\ref{ass:sparsity}, with $\cT_{\bar L}$ fixed and
finite,
\[
 \max_{(\tau_e,\tau_\mu)\in\cS_e\times\cS_\mu}
 \Bigl\lvert\sqrt n\{\thetahat(\tau_e,\tau_\mu)-\theta_0\}
  -\sqrt n\,\Pn\psi_0\Bigr\rvert=o_p(1),\qquad
 \psi_0=\psi(\cdot;\theta_0,\eta_0).
\]
\end{lemma}

\begin{remark}[This lemma closes the double use of $Y$; revision~1 lacked it]
\label{rem:uniform-why}
\textbf{This was the blocking gap in revision~1.} Lemmas
\ref{lem:invariance}, \ref{lem:selection}, \ref{lem:findim} do not compose into
Theorem~\ref{thm:main}: (i) Lemma~\ref{lem:invariance} is invariance of the
\emph{estimand}, i.e.\ of first-order bias, whereas the CLT needs invariance of
the \emph{influence function}; (ii) a finite mixture
$\sum_\tau\mathbf1\{\that=\tau\}\sqrt n\{\thetahat(\tau)-\theta_0\}$ is
\textbf{not} asymptotically normal merely because each component is, since the
selection indicators correlate with the components --- \emph{finiteness alone is
not enough}, what is needed is asymptotic equivalence to a \emph{fixed}
statistic; (iii) conditioning on $\that_\mu$ is not a repair, because leaf
membership is $X$-determined but \emph{which cells are merged} is $Y$-determined,
so $\mhat_\ell$ is conditionally biased given the selection event even on
$\{\that_\mu\in\cS_\mu\}$.

Lemma~\ref{lem:uniform} never conditions on a data-dependent event: it bounds
uniformly over a family not depending on the data, and a maximum of finitely many
$o_p(n^{-1/2})$ terms is $o_p(n^{-1/2})$. In the fixed tier this is free.

Worth stating alongside, because it is free and sharpens the novelty claim:
\textbf{neither tree uses treated outcomes.} $\that_e,\ehat$ are
$\sigma(X,A)$-measurable and $\that_\mu,\mhat$ are
$\sigma(X,A,\{Y_i:A_i=0\})$-measurable, so the only genuine double use of $Y$ is
of the \emph{control} arm; the double use of $A$ is harmless because the
propensity error enters \eqref{eq:skeleton} only multiplied by a residual whose
within-leaf control sum is exactly zero.
\end{remark}

\begin{lemma}[Conditional finite-dimensionality]\label{lem:findim}
For fixed $\tau$, the refitted nuisance is determined by at most $\bar L$ numbers
per nuisance, each a within-leaf mean converging at $O_p(n^{-1/2})$.
\end{lemma}

\begin{theorem}[Full-sample two-tree CLT]\label{thm:main}
Under A\ref{ass:causal}--A\ref{ass:sparsity}, with $\that_e,\that_\mu$ selected by
\eqref{eq:select} on all $n$ and leaf values refit on all $n$,
$\sqrt n(\thetahat-\theta_0)\rightsquigarrow N(0,V)$ where
\begin{equation}\label{eq:vdecomp}
  V=\pi^{-2}\E\Bigl[\ez\sigma_1^2+\frac{\ez^2\sigma_0^2}{1-\ez}
   +\ez(m_1-\mz-\theta_0)^2\Bigr]
  =\pi^{-2}\E\bigl[\{A(Y-\mz-\theta_0)-(1-A)\wz(Y-\mz)\}^2\bigr].
\end{equation}
\end{theorem}

\noindent \emph{Route:} \eqref{eq:skeleton} for each fixed pair, then
Lemma~\ref{lem:uniform}, then Lemma~\ref{lem:selection} to discard
$\{\that\notin\cS\}$, then $\pihat\to_p\pi$ and Remark~\ref{rem:sign}.

\begin{corollary}[Variance estimation and confidence intervals]\label{cor:variance}
With
\begin{equation}\label{eq:vhat}
 \Vhat=\frac1{n\pihat^2}\sum_{i=1}^n\Bigl[A_i\{Y_i-\mhat(X_i)-\thetahat\}
  -(1-A_i)\what(X_i)\{Y_i-\mhat(X_i)\}\Bigr]^2,
\end{equation}
$\Vhat\to_pV$ and $\thetahat\pm z_{1-\alpha/2}\sqrt{\Vhat/n}$ has asymptotic
coverage $1-\alpha$ --- pointwise in $P$, or uniformly over $\cP(\Delta_{\min})$
under A\ref{ass:margin-floor}.
\end{corollary}

\begin{remark}[Why $\Vhat$ is consistent, and why the obvious argument is wrong]
\label{rem:vhat}
\textbf{Corrected from revision~1}, which argued ``conditional on a sufficient
partition the problem is finite-dimensional''. Invalid twice: Lemma~\ref{lem:findim}
is a fixed-$\tau$ statement saying nothing about $\that$; and conditioning on
$\{\that=\tau\}$ conditions on a data-dependent event, under which the sample is
not i.i.d.\ from $P$, so no fixed-$\tau$ limit theorem applies.

The correct argument replaces conditioning by a \emph{union bound over the finite
sufficient class} and needs three inputs: Lemma~\ref{lem:invariance}, so every
$\tau\in\cS_j$ has the same limit (else $\operatorname{plim}\Vhat$ is a
selection-weighted mixture of different numbers); Lemma~\ref{lem:selection}, to
discard the non-sufficient event; and Lemma~\ref{lem:findim}, for fixed-$\tau$
consistency. Then $\lVert\widehat\eta-\eta_0\rVert_\infty=o_p(1)$ and
\eqref{eq:vhat} is a plug-in with $o_p(1)$ remainder by Cauchy--Schwarz.

Two things to advertise. $\Vhat$'s consistency needs only
$\lVert\widehat\eta-\eta_0\rVert_\infty=o_p(1)$ --- \emph{not} the rate condition
and \emph{not} Lemma~\ref{lem:cross}; since that lemma is the harder argument,
this isolates the risk. And there is no tension with invariance licensing only
distributional claims: $V=\E[\psi_0^2]$ is a second moment of a \emph{fixed
population} function, a functional of $P$ alone containing no selection, and
consistency is a probability limit, not a statement about
$\E\lVert\mhat-\mz\rVert^2$.

Finite-sample notes. By construction the summands of \eqref{eq:vhat} have exactly
zero sample mean, so it is simultaneously the uncentred second moment and the
centred sample variance --- no centring ambiguity. But because structure
\emph{and} leaf values are fit on all $n$, the summands are in-sample residuals,
so \eqref{eq:vhat} is downward biased by a relative factor
$\approx(2\bar L+1)/n$: asymptotically irrelevant, but exactly the mild
under-coverage that will appear in a simulation table and be misattributed to the
absence of splitting. A degrees-of-freedom divisor
$n-(2\card{\that_e}+\card{\that_\mu})$ is a sensible finite-sample convention. All
statements are on $\{\pihat>0\}$, whose probability tends to one.
\end{remark}

\begin{remark}[What $V$ is, and what efficiency may honestly be claimed]
\label{rem:efficiency}
\textbf{Substantially revised.} Revision~1 asserted that with $M$ fixed ``$P$ is
determined by finitely many parameters'', hence that this is not a nonparametric
efficiency result. \textbf{The premise is false}: with $M$ fixed, $P$ is
determined by $\{p_x,\ez(x)\}$ \emph{together with the conditional laws of $Y$
given $(A,X)$}, which A\ref{ass:moments} leaves unrestricted. The tangent space is
infinite-dimensional and $V$ is a genuine semiparametric bound; revision~1
conceded the paper's strongest interpretive claim on a false premise.

What is true: the ATT is a smooth function of finitely many identified
functionals of $P$ ($p_x,\ez,m_1,\mz$), and \emph{given a partition} the nuisance
problem is finite-dimensional in at most $2\bar L+1$ parameters --- which is why
no splitting is required --- while $P$ itself remains nonparametric. Attaining $V$
is therefore attaining the nonparametric ATT efficiency bound of
\citet{hahn1998}.

Efficiency \emph{within} the sparsity-restricted model is a different question,
and the answer is negative but interesting. Restricting $\ez$ to be constant on a
partition removes tangent directions, giving
\[
 V_\tau=V-\pi^{-2}\E\bigl[\ez(1-\ez)\{\Delta m-\Pi^v_\tau\Delta m\}^2\bigr]
  -(\text{an analogous }\mu\text{ term})\;\le\;V,
\]
with $\Delta m=m_1-\mz-\theta_0$ and $v_x=p_x\ez(x)\{1-\ez(x)\}$, strictly less
than $V$ whenever the treated--control contrast varies within the leaves of the
propensity partition. So efficiency \emph{is} left on the table --- Hahn's own
asymmetry, that for the ATT knowledge of the propensity lowers the bound. But
$V_{\mathrm{sparse}}=\max_\tau V_\tau$ depends \textbf{discontinuously} on which
partition is minimal, so any estimator attaining it must be irregular over the
union of sparse models (a Hodges/Leeb--P\"otscher construction: superefficient at
$P_0$, badly behaved in a shrinking neighbourhood). Declining to chase it is
therefore the right call, not a deficiency --- a far stronger statement than
conceding finite-dimensionality. Local asymptotic efficiency additionally needs a
regularity lemma along $\sqrt n$-local perturbations, which leave the sparse
class; plausible, probably cheap, not yet proved, so not claimed
(Section~\ref{sec:open}).
\end{remark}

\begin{remark}[Why two partitions, and why a shared one degenerates]
\label{rem:shared}
On $\that_e\wedge\that_\mu$ both $\what$ and $\mhat$ are leaf-constant, so
$\sum_{i\in\ell,A_i=0}(Y_i-\mhat_\ell)=0$ makes the entire control-arm term of
\eqref{eq:att-estimator} \emph{identically} zero and
$\thetahat_{\mathrm{int}}=\sum_\ell(n_{1,\ell}/n_1)(\bar Y_{1,\ell}-\bar
Y_{0,\ell})$, the treated-weighted stratified difference in means --- an exact
finite-sample identity, unaffected by clipping.

\textbf{But revision~1's gloss ``the propensity plays no role whatsoever''
overstates it}, and contradicts its own next clause: the propensity's fitted
\emph{values} are annihilated; its \emph{structure} still selects the
stratification and so changes both the point estimate and the variance. Since
$\thetahat_{\mathrm{int}}$ is consistent, asymptotically normal, and has $T=0$
identically (a much easier proof), a referee will ask why it is not recommended.
The defensible reasons are \emph{not} validity: (i) one displayed tree would carry
no fitted values, defeating the paper's purpose; and (ii)
$\card{\that_e\wedge\that_\mu}\le L_eL_\mu$, so the intersection can violate
A\ref{ass:mass}--A\ref{ass:budget} even when each tree satisfies them, inflating
$\sum_\ell\sigma^2_\ell/p_\ell$.
\end{remark}

\section{Part III: honest inference when sparsity fails}

Drop A\ref{ass:sparsity}; then Lemma~\ref{lem:invariance} fails and
Lemma~\ref{lem:biasbound} is the object to control.

\begin{definition}[Realised errors]\label{def:realised}
$\Dw=\Lw{\what-\wz}$ and $\Dm=\Lw{\mhat-\mz}$, at the fitted nuisances actually
reported.
\end{definition}

\begin{definition}[Approximation error, for comparison only]\label{def:delta}
$\delta_e=\min_{\tau\in\cT_{\bar L}}\Ltwo{\Pi_\tau\ez-\ez}$,
$\delta_\mu=\min_{\tau\in\cT_{\bar L}}\Lw{\Pi_\tau\mz-\mz}$.
\end{definition}

\begin{remark}[Why Part~III is on realised errors: the repair was mandatory]
\label{rem:whyD}
\textbf{The largest change from revision~1}, which used an allowance
$\delta_e\delta_\mu/(\pi c)$. Three defects, all fixed by
Definition~\ref{def:realised}.

\emph{(a) It was false.} \eqref{eq:bilinear} is exact in the \emph{fitted}
nuisances and $D_j=\delta_j+O_p(\sqrt{\bar L/n})$, so
\[
 \abs R\lesssim\frac1{\pi c^{3/2}}
  \Bigl[\delta_e\delta_\mu+(\delta_e+\delta_\mu)\sqrt{\bar L/n}+\bar L/n\Bigr],
\]
and revision~1 kept only the first term. With $\delta_e=n^{-0.05}$,
$\delta_\mu=n^{-0.50}$, $\bar L=n^{0.40}$ the product condition holds
($\delta_e\delta_\mu=n^{-0.55}$) while the omitted cross term is
$n^{-0.35}\gg n^{-1/2}$: coverage tends to $0$.

\emph{(b) The minimal patch reinstates what Part~I removes.} Repairing (a) by
assuming $(\delta_e\vee\delta_\mu)=o(\bar L^{-1/2})$ becomes, at
$\bar L\asymp\sqrt n$, exactly $\delta_j=o(n^{-1/4})$ \textbf{for each nuisance
separately} --- the per-nuisance condition Remark~\ref{rem:product} calls
unusable. Part~III would silently reinstate Part~I's headline correction.

\emph{(c) It bounded the wrong tree.} Without A\ref{ass:sparsity} there is no
margin (Lemmas \ref{lem:margin}, \ref{lem:selection} both assume it), so nothing
guarantees $\that_j$ is near-optimal in $\cT_{\bar L}$, and the realised error of
the \emph{reported} tree may exceed $\delta_j$ by a constant factor. Since
\eqref{eq:bilinear} is driven by the reported tree, $\delta_j$ is not the right
object.

Working with $\Dw,\Dm$ resolves all three, makes the constant exactly $1/\pi$,
removes $\bar L$ from the interval, and makes the allowance \emph{measurable}.
\end{remark}

\begin{remark}[Widen, do not subtract]\label{rem:no-subtract}
Trivially $\thetahat-\delta_{\mathrm{fid}}=\thetahat_{\mathrm{cf}}$, so
``debiasing'' by the fidelity diagnostic returns the cross-fitted estimator
exactly and forfeits the displayed-tree estimate for nothing. (Revision~1 stated
this as a numbered lemma; it is $a-(a-b)=b$, and is demoted here.) The substantive
point is that $\delta_{\mathrm{fid}}$ cannot serve as the widening term either:
under A\ref{ass:rate} and Theorem~\ref{thm:crossfit},
$\delta_{\mathrm{fid}}=(\thetahat-\theta_0)-O_p(n^{-1/2})$, so it is
$O_p(n^{-1/2})$ even when the bias is exactly zero and
$\PP(\abs{\delta_{\mathrm{fid}}}\ge\abs{\text{bias}})\not\to1$: neither centred on
nor dominating the bias.
\end{remark}

\begin{lemma}[The fidelity diagnostic is a specification test, not a correction]
\label{lem:spectest}
Under A\ref{ass:causal}--A\ref{ass:sparsity} and A\ref{ass:rate},
$\sqrt n\,\delta_{\mathrm{fid}}\rightsquigarrow N(0,V_{\mathrm{fid}})$ with
$V_{\mathrm{fid}}$ computable from the fold structure. Hence
$\delta_{\mathrm{fid}}$, compared with its own standard error, is a valid test of
A\ref{ass:sparsity}.
\end{lemma}

\begin{theorem}[Honest interval outside the class]\label{thm:honest}
Let $U_w,U_\mu$ be one-sided upper confidence bounds for $\Dw,\Dm$, each at level
$1-\alpha/4$, and let $\Vhat^{+}$ satisfy $\operatorname{plim}\Vhat^{+}\ge V$.
Then
\[
 \thetahat\;\pm\;\Bigl(z_{1-\alpha/2}\sqrt{\Vhat^{+}/n}
  +\frac{U_wU_\mu}{\pi}\Bigr)
\]
has asymptotic coverage at least $1-\alpha$ for $\theta_0$.
\end{theorem}

\begin{remark}[The plug-in $\Vhat$ may not be used here]\label{rem:vhat-anticons}
Without A\ref{ass:sparsity}, $\ehat\not\to\ez$ and $\mhat\not\to\mz$, so
\eqref{eq:vhat} converges elsewhere and $\operatorname{plim}\Vhat\ge V$ is
\emph{not} generally true. With $a=\mu_*-\mz$, $b=w_*-\wz$,
\[
 \operatorname{plim}\Vhat=V+2\pi^{-2}
  \bigl\{\E[(1-\ez)\wz b\sigma_0^2]-\E[\ez a\Delta m]\bigr\}
  +O(\lVert a\rVert^2+\lVert b\rVert^2),
\]
not sign-definite: $b\equiv0$ with $a=\epsilon(1-\ez)\Delta m$ gives a negative
first-order term against a second-order $O(\epsilon^2)$, so
$\operatorname{plim}\Vhat<V$. Theorem~\ref{thm:honest} therefore requires a
conservative $\Vhat^{+}$, whose construction is open (Section~\ref{sec:open}) ---
\emph{not} something the bias allowance repairs, since it addresses a different
term.
\end{remark}

\begin{proposition}[The realised errors are estimable when $M\ll n$]
\label{prop:estimable}
Suppose every cell is occupied in the relevant arm: $np_{\min}\gtrsim\log M$ for
$e$, and $ncp_{\min}\gtrsim\log M$ for $\mu$. With $\mhat_{\mathrm{sat}}$,
$\what_{\mathrm{sat}}$ the saturated cell-wise fits, the
degrees-of-freedom-corrected statistic
\[
 \widehat D_\mu^2=\sum_x\widehat p_x\{1-\ehat(x)\}
   \{\mhat(x)-\mhat_{\mathrm{sat}}(x)\}^2
  -\sum_x\widehat p_x\{1-\ehat(x)\}\frac{\widehat\sigma^2(x)}{n_{0,x}},
\]
and its analogue $\widehat D_w^2$, satisfy
$\widehat D_j^2-D_j^2=O_p(M/n+D_jn^{-1/2})$ up to logarithmic factors. Setting
$U_j^2=(\widehat D_j^2)_{+}+\kappa_na_{n,j}$ with $\kappa_n\to\infty$ slowly gives
valid one-sided bounds for Theorem~\ref{thm:honest}.
\end{proposition}

\begin{remark}[Three corrections that make Proposition~\ref{prop:estimable} honest]
\label{rem:prop-fixes}
\emph{(i) The plug-in is downward biased, so it cannot be used raw.} For squared
loss with an independent hold-out,
$\E[\widehat D^2_{j,\mathrm{plug}}]=D_j^2-(M-\bar L)\sigma^2/n+o(\cdot)$, negative
because $M\gg\bar L$, plus a same-signed winner's-curse term from the
minimisation. At $D_j^2=2M\sigma^2/n$ the plug-in returns about $0.71D_j$ --- a
$29\%$ understatement per factor, roughly $50\%$ on the product. That is a
coverage failure of the same character as Remark~\ref{rem:vhat-anticons}'s, and it
compounds with it. Hence the explicit correction above.

\emph{(ii) An upper bound must be substituted, not a point estimate.} Note the
non-smoothness of $\sqrt\cdot$ at $0$ does \emph{not} threaten honesty, since
$\sqrt\cdot$ is monotone and a bound on $D^2$ maps to one on $D$; it costs
\emph{width} (Corollary~\ref{cor:width}).

\emph{(iii) The two one-sided statements must be composed}, requiring
$\PP(U_w\ge\Dw\text{ and }U_\mu\ge\Dm)\to1$: hence $\alpha/4$ per side.

Two further notes. The rate $M/n+D_jn^{-1/2}$ is sharper than the naive
$\sqrt{M/n}$, and the sharpening is load-bearing: with $\sqrt{M/n}$ the width
exceeds $n^{-1/2}$ whenever $M\to\infty$ at \emph{any} rate. And the weighted norm
pays off again --- the aggregate degrees-of-freedom term is governed by $n$ alone,
since $\sum_xp_x(1-\ez)\sigma^2(x)/\{np_x(1-\ez)\}=n^{-1}\sum_x\sigma^2(x)$, with
no $1/c$.
\end{remark}

\begin{corollary}[Width, and when the price is negligible]\label{cor:width}
The half-width of Theorem~\ref{thm:honest} is
\[
 \asymp\max\bigl\{n^{-1/2},\ \Dw\Dm,\ (M/n)^{1/2}(\Dw\vee\Dm),\ M/n\bigr\},
\]
collapsing to $\max\{n^{-1/2},\Dw\Dm\}$ if and only if $M=o(\sqrt n)$ and
$(\Dw\vee\Dm)=o(M^{-1/2})$.
\end{corollary}

\begin{remark}[Scope: two constraints, and revision~1 reported only the permissive
one]\label{rem:scope}
Estimating $D_j$ needs the saturated fit to \emph{exist} --- the occupancy
condition, roughly $30$ observations per cell, i.e.\ $M\lesssim n/30$. Keeping the
width at $\max\{n^{-1/2},\Dw\Dm\}$ needs $M=o(\sqrt n)$, i.e.\
$d\lesssim\tfrac12\log_2n$ for binary coordinates. \textbf{These differ by a
square root} and revision~1 tabulated only the first:

\begin{center}\begin{tabular}{@{}rrrr@{}}
\toprule
$n$ & $M$ for existence & $M$ for cheap width & binary $d$ (cheap) \\ \midrule
$1{,}000$ & $\lesssim33$ & $\lesssim31$ & $d\le4$ \\
$10{,}000$ & $\lesssim333$ & $\lesssim100$ & $d\le6$ \\
$100{,}000$ & $\lesssim3{,}333$ & $\lesssim316$ & $d\le8$ \\
\bottomrule
\end{tabular}\end{center}

So the honest figures are $d\le6$ at $n=10^4$ and $d\le8$ at $n=10^5$, not $7$--$8$
and $10$--$11$. The binding constraint is the \emph{smallest} cell, not the
average; and the control share is $\E[1-\ez]\ge c$, which under one-sided
positivity may be far below $1/2$, so revision~1's ``effective sample is about
$n/2$'' is replaced by $n\min_xp_x\{1-\ez(x)\}$. This constrains
Proposition~\ref{prop:estimable} only: Theorems \ref{thm:main} and
\ref{thm:honest} leave $d$ unconstrained, where the binding constraint is
interpretability itself.
\end{remark}

\begin{remark}[The assumed-decay alternative, and an honest comparison]
\label{rem:assumed-decay}
The alternative to measuring is assuming: $\delta_j(L)\le C_jL^{-q_j}$ with
$q_e+q_\mu>1$ --- the discrete analogue of assuming a H\"older or Besov ball via
Jackson-type theorems, and nobody objects to H\"older. Supply a sufficient
condition as a lemma: weak-$\ell_r$ decay of tree/Haar coefficients gives
$q=1/r-1/2$. No decay rate is implied by finiteness: parity on $d$ binary
coordinates has $\delta(L)$ bounded away from zero for every axis-aligned tree
with $L<2^{d-1}$ leaves and then drops to zero, so decay can be a step function of
$L$. ANOVA or low-order-interaction decay must \textbf{not} be used --- such
functions are well approximated by \emph{additive} models, trees' worst case
(Remark~\ref{rem:additive}).

Honest comparison, which revision~1 elided by claiming
Proposition~\ref{prop:estimable} needs ``no assumed decay rate'': it needs
occupancy, $M=o(\sqrt n)$, and $(\Dw\vee\Dm)=o(M^{-1/2})$, the last two no more
checkable than the decay assumption. The defensible claim is narrower --- \emph{the
allowance can be replaced by a data-driven upper bound, so the interval is honest
without knowing $q_e,q_\mu$; the price is the extra
$(M/n)^{1/2}(\Dw\vee\Dm)+M/n$ width}. The assumed-decay route yields the strictly
narrower $\max\{n^{-1/2},\delta_e\delta_\mu\}$ and remains the main line.

Finally, with $\bar L$ fixed the $D_j$ do not shrink, so the price does not
vanish. That is Part~III's honest message: interpretability carries a cost that is
measurable and, unless the truth is structurally sparse, does not disappear.
Displaying the held-out $\widehat D_j(L)$ path along the optimal-tree
$L$-sequence is the right way to show it.
\end{remark}

\section{Changes from revision 1}\label{sec:changes}

Four blocking defects and their repairs:
\begin{enumerate}[label=\textbf{B\arabic*.},leftmargin=*,nosep]
  \item \textbf{The honest interval was false} --- it bounded approximation rather
  than realised error, omitting a cross term that can dominate, and the minimal
  patch would have reinstated Part~I's per-nuisance $n^{-1/4}$ condition. Repaired
  by Def.~\ref{def:realised}, Lemma~\ref{lem:biasbound},
  Theorem~\ref{thm:honest} (Remark~\ref{rem:whyD}).
  \item \textbf{``Refines $\tau^*_j$'' broke Theorem~\ref{thm:main}} --- a minimal
  box partition need not be tree-realisable; the pinwheel gives
  $\cS_j=\emptyset$. Repaired by Def.~\ref{def:sufficient}
  (Remark~\ref{rem:pinwheel}).
  \item \textbf{Theorem~\ref{thm:main}'s proof did not compose} --- invariance of
  the estimand is not invariance of the influence function, and a finite mixture
  is not asymptotically normal componentwise. Repaired by
  Lemma~\ref{lem:uniform}, which also repairs Remark~\ref{rem:vhat}.
  \item \textbf{$\mhat$ could be undefined} --- A\ref{ass:mass} counted
  observations, not controls. Repaired with a control count and a fallback
  convention.
\end{enumerate}

Further corrections: positivity weakened to one-sided (A\ref{ass:causal}); the
bias constant is exactly $1/\pi$, and the norm-equivalence constant for norms is
$c^{-1/2}$ not $1/c$ --- the latter caused the former
(Remark~\ref{rem:normconst}); $V>0$ added; the leaf budget pinned to fixed with
the growing regime scoped out (Remark~\ref{rem:fixed-tier});
Lemma~\ref{lem:cross}'s $\log n$ removed and the rate sharpened to
$O_p(\sqrt{L_e}/n)$; the claim that clipping requires $m_n\to\infty$ withdrawn,
and the clip level tied to $c$ (Remark~\ref{rem:clip}); the efficiency framing
reversed to reclaim the nonparametric bound (Remark~\ref{rem:efficiency});
coverage declared pointwise with an optional uniformity class
(A\ref{ass:margin-floor}); the plug-in $\Vhat$ shown inadmissible in
Theorem~\ref{thm:honest} (Remark~\ref{rem:vhat-anticons});
Proposition~\ref{prop:estimable} given a degrees-of-freedom correction, one-sided
bounds and $\alpha$-splitting; the width order corrected
(Corollary~\ref{cor:width}) and the scope table recomputed
(Remark~\ref{rem:scope}); the vacuous no-subtract lemma demoted and replaced by
Lemma~\ref{lem:spectest}; $2^s$ corrected to $J^s$ (Remark~\ref{rem:additive});
Remark~\ref{rem:shared}'s ``no role whatsoever'' withdrawn; the $\pihat$ display
corrected (Remark~\ref{rem:sign}); the control-arm cancellation in the weighted
norm made explicit (Remark~\ref{rem:mass}).

\section{Open items}\label{sec:open}

\begin{enumerate}[nosep,label=(\arabic*)]
  \item \textbf{A conservative variance estimator for Part~III}: $\Vhat^{+}$ with
  $\operatorname{plim}\Vhat^{+}\ge V$; the plug-in does not qualify.
  \item \textbf{Proofs}: Lemma~\ref{lem:uniform} (assembly),
  Lemma~\ref{lem:cross} in full, $V_{\mathrm{fid}}$ in
  Lemma~\ref{lem:spectest}, Theorem~\ref{thm:honest}, and
  Proposition~\ref{prop:estimable}'s rate.
  \item \textbf{Local regularity} along $\sqrt n$-local perturbations, needed
  before claiming local asymptotic efficiency.
  \item \textbf{The growing regime} ($M_n\to\infty$ or $L_n\to\infty$): entropy
  $\sqrt{L_n\log(dJ)}$, a bounded or sub-Gaussian outcome, possibly splitting for
  structure selection. Deliberately untreated.
  \item \textbf{Part~I's oracle inequality} bounds terms by
  $O_p(\sqrt{L_n\log n/n})$ and then squares them inside an additive risk
  decomposition, understating the rate by a square; it needs the localised
  sub-root argument.
\end{enumerate}

\begin{thebibliography}{9}
\bibitem[Chen et al.(2022)]{chen2022} Chen, Q., Syrgkanis, V., Austern, M.
  (2022). Debiased machine learning without sample splitting. arXiv:2206.01825.
\bibitem[Hahn(1998)]{hahn1998} Hahn, J. (1998). On the role of the propensity
  score in efficient semiparametric estimation of average treatment effects.
  \emph{Econometrica} 66, 315--331.
\bibitem[Kennedy(2024)]{kennedy2024} Kennedy, E. H. (2024). Semiparametric
  doubly robust targeted double machine learning: a review.
\bibitem[Robins \& Rotnitzky(1995)]{robins1995} Robins, J. M., Rotnitzky, A.
  (1995). Semiparametric efficiency in multivariate regression models with
  missing data. \emph{JASA} 90, 122--129.
\bibitem[van der Laan \& Robins(2003)]{vanderlaan2003} van der Laan, M. J.,
  Robins, J. M. (2003). \emph{Unified Methods for Censored Longitudinal Data and
  Causality}. Springer.
\end{thebibliography}

\end{document}
